\documentclass{article}

\usepackage{amsmath, amsthm, amssymb}
\usepackage{graphicx}
\usepackage{verbatim}
\usepackage{natbib}
\usepackage{caption}
\usepackage{subcaption}
\usepackage{fancyvrb}
\usepackage{enumerate}
\usepackage{relsize}
\usepackage{multirow}
\usepackage{mathrsfs}

\usepackage{hyperref}
\usepackage[margin=1.4in]{geometry}
\hypersetup{colorlinks,citecolor=blue,urlcolor=blue,linkcolor=blue}
\newcommand{\Lihua}[1]{{\color{blue}#1}}

\newcommand\blfootnote[1]{%
  \begingroup
  \renewcommand\thefootnote{}\footnote{#1}%
  \addtocounter{footnote}{-1}%
  \endgroup
}

\usepackage{stefan_tex}
\graphicspath{{./figures/}}

%\numberwithin{equation}{section}

%%%%%%%%% Theorems

\theoremstyle{definition}
\newtheorem{exam}{Example}
\newtheorem{defi}{Definition}
\newtheorem{assumption}{Assumption}
\newtheorem{property}{Property}

 \theoremstyle{plain}
 \newtheorem{prop}{Proposition}
 \newtheorem{conj}[prop]{Conjecture}
 \newtheorem{coro}[prop]{Corollary}
 \newtheorem{lemm}[prop]{Lemma}
 \newtheorem{theo}[prop]{Theorem}
%\theoremstyle{plain}
%\newtheorem{prop}{Proposition}[section]
%\newtheorem{coro}{Corollary}[section]
%\newtheorem{lemm}{Lemma}[section]
%\newtheorem{theo}{Theorem}[section]

\theoremstyle{remark}
\newtheorem{comm}{Comment}
\newtheorem{rema}{Remark}

\author{Roshni Sahoo\\
  \texttt{r.sahoo@columbia.edu}}


\title{Related Work Data Fusion}

\date{Columbia University}


\begin{document}

\maketitle

\begin{section}{Regression Models and Life Tables}
We consider a population of individuals; for each individual, we observe either the time to ``failure'' or the time to ``loss'' or censoring. For censored individuals, we know only that the time to failure is greater than the censoring time.

Let $T$ be a random variable representing failure time. Let $F(t) = \PP[]{T \geq t}$. Let $\lambda(t)$ be the hazard, or age-spepcific failure rate, i.e. 
\[\lambda(t) = \lim_{\Delta \rightarrow 0^{+}} \frac{\PP[]{t \leq T \leq t + \Delta \mid T \geq t}}{\Delta}.\]
The probability that failure happens at time $t$, given that it has not yet occured. If $T$ is discrete, then 
\[ \lambda(t) = \sum \lambda_{u_{j}} \delta(t - u_{j}),\]
where $\lambda_{t} = \PP[]{T=t \mid T \geq t}$.

Let $n$ individuals be observed to failure and the rest be censored. The only information about the failure time of a censored individual is that it exceeds the censoring time. The failure times are given by 
\[ t_{(1)} < t_{(2)} < \dots < t_{(n)}.\]
Let $m_{(i)}$ be the number of units with failure time equal to $t_{(i)}$. The risk set $\mathcal{R}(t)$ is the set of individuals at risk at time $t$. Let $r_{(i)}$ be the number of such individuals for $t=t_{(i)}$. The product limit estimate of the distribution is obtained by taking estimated conditional probabilities that agree with the observed distributional frequencies:
\[ \hat{\lambda}(t) = \sum_{i=1}^{k} \frac{m_{(i)}}{r_{(i)}} \cdot \delta(t - t_{(i)}).\]
MLE estimate in the family of all possible distributions.

Suppose we also observe measurements $x$ (covariates) for each unit in addition to failure time. Goal is to assess relationship between failure times and $x$. Posit a model where the hazard is 
\[ \lambda(t; x) = \exp(z_{(i)} \beta) \cdot \lambda_{0}(t).\]
The probability of failure at time $t_{(i)}$ is given by
\[ \exp(z_{(i)} \beta) \Big/ \sum_{l \in \mathcal{R}(t_{(i)})} \exp(z_{(l)} \beta).\]


Our stationarity assumption implies that time-dependent factors have only a multiplicative impact on the probability of response.

\end{section}

\end{document}